\documentclass[10pt,a4paper]{amsart}
\usepackage[margin=19mm]{geometry}
\usepackage{amsmath,amssymb,mathtools}
\usepackage[scr=rsfs,cal=euler]{mathalfa}
\usepackage{xfrac}
\usepackage[unicode,hidelinks,bookmarks=false]{hyperref}
\newcommand{\ie}{i.e.\ }
\newcommand{\cf}{cf.\ }
\newcommand{\resp}{respectively\ }
\newcommand{\Subsection}{\subsection}
\providecommand{\nobreakdash}{\nobreak\hbox{-}}
\setlength{\emergencystretch}{2em}
\multlinegap0pt
\usepackage{bm}
\usepackage{bbm}
\usepackage{dsfont}
\usepackage{xspace}
\usepackage{cancel}
\usepackage{mathtools}
\usepackage{amsthm}
\makeatletter
\@ifundefined{theorem}{\newtheorem{theorem}{Theorem}}{}
\@ifundefined{lemma}{\newtheorem{lemma}[theorem]{Lemma}}{}
\@ifundefined{proposition}{\newtheorem{proposition}[theorem]{Proposition}}{}
\makeatother
\title[Single Lie Brackets on Closed Manifolds]{Single Lie Brackets on Closed Manifolds}
\author{Zeyu Zeng}
\date{Source: arXiv:2609.05958v2 (CC BY 4.0)}
\begin{document}
\maketitle

\noindent\emph{Source and licence.} Single Lie Brackets on Closed Manifolds. Version arXiv:2609.05958v2 (CC BY 4.0), distributed under the Creative Commons Attribution 4.0 International licence, as recorded in the arXiv licence metadata for this exact version and shown on the abstract page https://arxiv.org/abs/2609.05958v2. Full rights notice: licenses/arxiv-2609.05958-CC-BY-4.0.txt.\par\smallskip

\noindent\textit{Editorial scope.} Counted unit: the Proposition whose source label is \texttt{prop:graph} in the article (source0.tex lines 553--560, source label \texttt{prop:graph}), delivered together with the article's own complete original proof of it (source0.tex lines 562--762, one \texttt{proof} environment of 201 lines). No mathematics is written, repaired or re-derived: statement and proof are the article's own text line for line. Required context retained uncounted: lem:scalar (lines 280--282); lem:positive (lines 337--359). Every definition, display and label of those lines is retained unchanged. Omitted from this package and not redistributed: 1--279, 283--336, 360--552, 561--561, 763--896. Nothing inside a retained excerpt is omitted or altered: every retained body is delivered exactly as the article's lines are written, so no omission receipt of any kind accompanies this package. Citations: the retained excerpts carry no citation command at all, so no bibliography entry of the article is transcribed and no reference is redistributed. Reference adaptation: none was needed and none was made; every reference inside the delivered unit resolves inside it, so no unresolved reference or citation is produced. Printed slips kept literal: no printed word, symbol, hypothesis or number of the counted statement or of its proof was altered. Layout and class: the article's own class is \texttt{article} with packages including fontenc, lmodern, amsmath, amssymb, amsthm, mathtools, biblatex, geometry, microtype, hyperref, enumitem; the wrapper is the lane's pinned \texttt{amsart} editorial wrapper at 10pt on A4, which restates the theorem-like environments the retained text needs and adds the article's own macro definitions that the retained text needs (none), with extra wrapper packages (bm, bbm, dsfont, xspace, cancel, mathtools). The delivered numbering is the unit's own. Assets: no retained span contains a figure environment, a graphics-inclusion command, a PDF-page inclusion, a table or a TikZ picture; no image file is copied into this package, the package declares no asset dependency and no image file is redistributed with it. Licence: version arXiv:2609.05958v2 of this article is distributed under the Creative Commons Attribution 4.0 International licence, as recorded in the arXiv licence metadata for this exact version and shown on the abstract page https://arxiv.org/abs/2609.05958v2. A full rights notice is delivered at licenses/arxiv-2609.05958-CC-BY-4.0.txt. Characters: every retained excerpt is pure ASCII, so the delivered engine, XeLaTeX with its default Latin Modern text font, reports no missing character.
\begin{lemma}[Scalar connector]\label{lem:scalar}
Let $c_{\mathrm{in}},c_{\mathrm{out}}$ be smooth functions on a transverse domain in $\mathbb R^m$, $m\ge1$, whose difference has compact support in its interior. In any prescribed positive time interval there is a smooth pair with bracket $T$, entrance germ $R_{c_{\mathrm{in}}}$ and exit germ $R_{c_{\mathrm{out}}}$. It retains that common germ off a compact subset of the transverse domain.
\end{lemma}

\begin{lemma}[Positive normal component]
\label{lem:positive}
Fix $\varepsilon\in\{1,-1\}$. Let $E\subset C$ be closed and
suppose that, for some $R_0>0$,
$$
  \partial C\subset E,
  \qquad
  [0,L]\times\{u:|u|\geq R_0\}\subset E.
$$
Let $N$ be a relatively open neighborhood of $E$ in $C$.
Suppose that smooth vector fields $Y_*,Z_*$ are given on $N$ and
satisfy
$$
  [Y_*,Z_*]=\varepsilon S
  \qquad\text{on }N.
$$
Then there exist smooth vector fields $Y,Z$ on $C$,
agreeing with $Y_*,Z_*$ on a neighborhood of $E$, such that
$$
  \varepsilon\,ds([Y,Z])\geq\frac12
  \qquad\text{on }C.
$$
\end{lemma}

\begin{proposition}[Graph connector with every lateral side attached]
\label{prop:graph}
In coordinates $(t,u)$, let $X=\partial_t$ and $R=(X,tX)$. Let $P$ be a closed disk in the transverse coordinate domain, and let $\Gamma=\{t=g(u):u\in P\}$ be a smooth graph. Suppose an exact pair $(Y_0,Z_0)$ is given near $\Gamma$, with $[Y_0,Z_0]=X$, and equals $R$ on a neighborhood of its boundary. For either sign $\varepsilon$ and every sufficiently small $H>0$, there is an exact pair on a neighborhood of the graph cylinder
$$
 \{(t,u):u\in P,\quad t=g(u)+\varepsilon r,\quad0\le r\le H\}
$$
which equals the old pair near $r=0$, equals $R$ near $r=H$, and equals $R$ on a full vertical collar of $\partial P$.
\end{proposition}

\begin{proof}
Fix $\varepsilon\in\{1,-1\}$ and a sufficiently small $H>0$. Choose $L,\ell>0$ such that
$$
    2L+\ell<H.
$$
We shall first construct an exact pair up to a variable exit graph, correct its outgoing germ in a band of width $\ell$, and then extend by the reference pair.

Extend $g$ smoothly to $\mathbb R^m$, making it constant outside a compact set. Since the old pair agrees with $R$ near the boundary of the entrance graph, its entrance germ extends to all transverse variables by $R$. This extension is reference both outside $P$ and on a neighborhood of $\partial P$.

Consider the auxiliary cylinder $C_L=[0,L]\times\mathbb R^m$ with $(s,v)\in C_L$ and $S=\partial_s$ Under the coordinate map
$$
    J(s,v)=(g(v)+\varepsilon s,v),
$$
the target field and reference pair become
$$
    J^*X=\varepsilon S,
    \qquad
    J^*R=
    \bigl(\varepsilon S,\,[g(v)+\varepsilon s]\varepsilon S\bigr).
$$
Prescribe the pulled-back old germ near $s=0$, and prescribe $J^*R$ near $s=L$, throughout a full vertical collar of $\partial P$, and for all $v\notin P$. These prescriptions agree on their overlaps. After choosing disjoint endpoint collars, Lemma~\ref{lem:positive} gives a smooth pair $(\widehat Y,\widehat Z)$ retaining these germs
and satisfying
$$
    L_0=[\widehat Y,\widehat Z],
    \qquad
    ds(\varepsilon L_0)\ge\frac12.
$$
Set $V=\varepsilon L_0$. In every prescribed region, $V=S$.

We next use the flow of $V$ to straighten the bracket. Its coefficients are bounded on $C_L$: outside $P$ the field is $S$, while $[0,L]\times P$ is compact. The continuation theorem for smooth ODEs therefore rules out finite-time escape while a trajectory remains in the cylinder. Moreover, along every trajectory,
$$
    \frac{ds}{dr}=ds(V)\ge\frac12.
$$
Thus the trajectory starting at $(0,u)$ reaches $s=L$ exactly once, at a time $T(u)$ satisfying
$$
    0<T(u)\le 2L.
$$
The hit is transverse, so the implicit function theorem shows that $T$ is smooth.

Define
$$
    D_T=
    \{(r,u):u\in\mathbb R^m,\ 0\le r\le T(u)\},
    \qquad
    \Theta(r,u)=\operatorname{Fl}_V^{\,r}(0,u).
$$
This is a diffeomorphism from $D_T$ onto $C_L$. Indeed, the same continuation and monotonicity arguments show that every point of $C_L$ reaches $s=0$ uniquely under backward flow, in time at most $2L$. If $\tau(s,v)$ denotes this backward travel time and $\pi$
is the transverse projection, then
$$
    \Theta^{-1}(s,v)
    =
    \left(
        \tau(s,v),
        \pi\operatorname{Fl}_V^{-\tau(s,v)}(s,v)
    \right).
$$
The implicit function theorem at the bottom hit makes $\tau$, and hence this inverse, smooth. The endpoint germs allow these statements to be understood on neighborhoods of the closed cylinders.

By construction,
$$
    \Theta_*\partial_r=V=\varepsilon L_0,
    \qquad
    \Theta^*L_0=\varepsilon\partial_r.
$$
Now introduce the physical source coordinates through
$$
    F(r,u)=(g(u)+\varepsilon r,u).
$$
Since $F_*(\varepsilon\partial_r)=X$, the pair
$$
    (Y_1,Z_1)=F_*\Theta^*(\widehat Y,\widehat Z)
$$
satisfies
$$
    [Y_1,Z_1]
    =
    F_*\Theta^*L_0
    =
    X.
$$

We record the relative properties of this construction. Near the entrance, $V=S$, so $\Theta(r,u)=(r,u)$ for sufficiently small $r$. Consequently $(Y_1,Z_1)$ retains the old entrance germ. On the prescribed vertical collar, every integral curve of $V$ is vertical. By uniqueness, any trajectory meeting this collar must coincide with the vertical trajectory through that point; in particular, trajectories cannot cross $\partial P$. It follows that $\Theta$ restricts to a diffeomorphism
$$
    \{(r,u):u\in P,\ 0\le r\le T(u)\}
    \longrightarrow
    [0,L]\times P.
$$
On the collar, and also outside $P$, we have
$$
    \Theta(r,u)=(r,u),
    \qquad
    T(u)=L.
$$
Thus the physical pair is reference throughout the retained vertical collar.

It remains to identify and correct the exit germ. Define the transverse exit map $H_0$ by
$$
    \Theta(T(u),u)=(L,H_0(u)).
$$
Since $V=S$ near the auxiliary top, the flow has the explicit form
$$
    \Theta(r,u)=\bigl(L+r-T(u),H_0(u)\bigr)
$$
near $r=T(u)$. In particular, $\Theta^*S=\partial_r$ there.
Pulling back the auxiliary reference pair therefore gives
$$
    \Theta^*(\widehat Y,\widehat Z)
    =
    \left(
        \varepsilon\partial_r,\,
        \bigl[g(H_0(u))+\varepsilon(L+r-T(u))\bigr]
        \varepsilon\partial_r
    \right).
$$
Because $t=g(u)+\varepsilon r$, its physical expression is
$$
    (Y_1,Z_1)=\bigl(X,(t+c(u))X\bigr),
$$
where
\begin{equation}\label{eq:defect}\tag{4.3}
    c(u)=g(H_0(u))-g(u)+\varepsilon(L-T(u)).
\end{equation}
The function $c$ is smooth. On the retained collar and outside $P$, we have $H_0(u)=u$ and $T(u)=L$; hence $c=0$ there. In particular, $\operatorname{supp}c$ is a compact subset of $\operatorname{int}P$.

To apply Lemma~\ref{lem:scalar} above the variable exit graph, put
$$
    b(u)=g(u)+\varepsilon T(u),
    \qquad
    \rho=r-T(u),
$$
and use the coordinate map
$$
    K(\rho,u)=(b(u)+\varepsilon\rho,u).
$$
The graph $\rho=0$ is precisely the physical exit graph. In these coordinates,
$$
    K^*X=\varepsilon\partial_\rho,
$$
and the incoming pair is
$$
    K^*(Y_1,Z_1)
    =
    \left(
        \varepsilon\partial_\rho,\,
        [\rho+\varepsilon(b(u)+c(u))]\partial_\rho
    \right).
$$
Multiplying only the first field by $\varepsilon$ gives the normalized
pair
$$
    \left(
        \partial_\rho,\,
        [\rho+\varepsilon(b(u)+c(u))]\partial_\rho
    \right),
$$
whose bracket is $\partial_\rho$.

Apply Lemma~\ref{lem:scalar} in the band $0\le\rho\le\ell$, with offsets
$$
    c_{\mathrm{in}}=\varepsilon(b+c),
    \qquad
    c_{\mathrm{out}}=\varepsilon b.
$$
Their difference is compactly supported in $\operatorname{int}P$. We may therefore take the modification to have compact transverse support in $\operatorname{int}P$, retaining a smaller full collar of $\partial P$. Let $(A,B)$ be the resulting normalized pair. Thus
$$
    [A,B]=\partial_\rho,
$$
and its entrance and exit germs are the prescribed scalar pairs.

Undo the normalization and the coordinate change by setting
$$
    (Y_2,Z_2)=K_*(\varepsilon A,B).
$$
Then
$$
    [Y_2,Z_2]
    =
    K_*(\varepsilon\partial_\rho)
    =
    X.
$$
Its entrance germ agrees with $(Y_1,Z_1)$. At its exit,
$$
    (Y_2,Z_2)
    =
    K_*\left(
        \varepsilon\partial_\rho,\,
        [\rho+\varepsilon b(u)]\partial_\rho
    \right)
    =
    (X,tX)=R.
$$
On the retained vertical collar, the correction is also reference throughout.

The corrected upper graph is $r=T(u)+\ell$. Since
$$
    T(u)+\ell\le 2L+\ell<H,
$$
extend the pair by $R$ from this graph to $r=H$. Each gluing takes place where the adjoining pairs agree on an open
neighborhood, so the resulting pair is smooth and has bracket $X$. It retains the old germ near $r=0$, is reference near $r=H$, and is reference on a full vertical collar of $\partial P$. Finally, compactness of $P$ allows the endpoint neighborhoods to be chosen uniformly, giving the full collars asserted in the proposition.
\end{proof}

\end{document}
